\section{QARMAs}

\begin{enumerate}
  \item How can the WTO define a legitimate "emergency in international relations" under GATT Article XXI without infringing on states' rights to self-define their own security interests — and should this definition expand to cover modern threats (cyber, critical minerals, data, supply-chain resilience)?
  \item How can the WTO distinguish between legitimate supply chain resilience and disguised mercantilism?
  \item What concrete institutional concession can be made to address the United States’ concerning “judicial overreach”, in order to secure the American unblocking of the Appellate Body judge appointments?
  \item Should the Multi-Party Interim Appeal Arbitration Arrangement (MPIA) be formalised as a permanent parallel system, or should all diplomatic effort be directed at restoring a universal Appellate Body?
  \item How can the WTO regulate or integrate exclusive economic alliances (such as AUKUS, technology clubs, or critical mineral partnerships) to ensure they do not permanently violate the Most-Favoured-Nation (MFN) principle?
  \item What legal disciplines should govern unilateral coercive measures, asset freezes, and export controls invoked on national security grounds, particularly regarding their extraterritorial impact on third-party states?
  \item If a state's measures are found by a WTO panel to violate Article XXI (i.e., not implemented during a legitimate international emergency), what remedies or non-escalatory countermeasures should the WTO authorise?
  \item How can the WTO coordinate with the UN to establish binding "humanitarian exemptions" within trade sanctions, ensuring critical food, agricultural inputs, and medical supply chains remain insulated from geoeconomic conflict?
  \item Who should have the authority to classify a good as “dual-use” for security purposes, and what disciplines, if any, should apply to prevent overclassification as a backdoor trade barrier?
  \item Should export controls on dual-use and advanced technologies be subject to WTO-level multilateral criteria (as under the Wassenaar Arrangement) or does the current fragmented, unilateral approach best reflect legitimate sovereign discretion?
  \item Do large-scale, security-justified industrial subsidies (e.g. the CHIPS Act, the Green Deal Industrial Plan) fall outside WTO subsidy disciplines when framed as national security measures, and should the SCM Agreement be updated to address this?
  \item What specific notification and consultation obligations should accompany a state's invocation of Article XXI, and what consequences (if any) should follow non-compliance?
\end{enumerate}
